\newpage

% =========================================================================
% CAPÍTULO I: GENERALIDADES
% =========================================================================
\chapter{Generalidades}

\section{Planteamiento del problema}

\subsection{Descripción del problema}
El desarrollo de herramientas computacionales para lenguas con menos recursos digitales enfrenta limitaciones de datos, representación y evaluación. Para el quechua, los trabajos disponibles documentan tanto la creación de corpus como la formalización de recursos léxicos, y también muestran que la variedad y el tipo de información lingüística afectan las tareas de traducción automática \parencite{Orcotoma2025,Duran2026,Court2024}. Estos antecedentes no permiten atribuir una única causa ni una magnitud uniforme a las dificultades de procesamiento del Quechua Collao.

Los modelos que dependen de servicios remotos no pueden atender solicitudes cuando el dispositivo carece de conexión. La ejecución local evita esa dependencia durante la inferencia, pero obliga a trabajar dentro de los límites de memoria y procesamiento disponibles; las técnicas de cuantización pueden reducir el tamaño del modelo a costa de compromisos que deben medirse en el equipo objetivo \parencite{Xu2024}.

La generación aumentada por recuperación (RAG) incorpora fragmentos de documentos al contexto de entrada para que la respuesta pueda apoyarse en información externa. Su desempeño depende de la recuperación y de la capacidad del modelo para usar el contexto. En una evaluación de tutoría educativa, un modelo pequeño tuvo dificultades con contextos extensos o irrelevantes; ese resultado evidencia un riesgo de diseño, no una conclusión directa sobre Quechua Collao \parencite{Hevia2025}.

El ajuste de bajo rango sobre modelos cuantizados ofrece una alternativa al entrenamiento completo, pero sus necesidades de hardware dependen del tamaño del modelo y de la configuración. El trabajo original de QLoRA informa resultados en equipos con recursos considerablemente mayores que la GPU local de 4 GB considerada en este proyecto, por lo que no se asumirá que el ajuste será viable antes de probarlo \parencite{Dettmers2023}. El problema de investigación consiste, por tanto, en determinar cómo se comporta una arquitectura compacta que combine recuperación local y adaptación bajo las restricciones reales del entorno.

\subsection{Identificación del problema}
El problema central es la falta de evidencia experimental sobre la calidad lingüística, la fidelidad al contexto recuperado y el costo computacional de un sistema local de generación en Quechua Collao que combine un modelo compacto y RAG, con adaptación QLoRA evaluada de manera condicionada a la memoria disponible.

La Figura~\ref{fig:arbol_problemas} sintetiza, mediante un árbol de problemas, las causas identificadas y los efectos derivados del problema central.

\begin{figure}[H]
\centering
\caption{Árbol de problemas: causas tecnológicas y efectos computacionales}
\label{fig:arbol_problemas}
\resizebox{\textwidth}{!}{%
\begin{tikzpicture}[every node/.style={text width=3.3cm, minimum height=1.7cm}]
\node[tcentral, text width=6.6cm, minimum height=1.8cm] (centro) at (0,0) {\textbf{Problema central}\\[2pt]Falta de evidencia sobre calidad y fidelidad de respuestas con RAG local para Quechua Collao};

\node[tnaranja] (c1) at (-5.5,-3.2) {\textbf{Causa 1: datos}\\[2pt]Disponibilidad y cobertura de fuentes documentales bilingües};
\node[tnaranja] (c2) at (0,-3.2) {\textbf{Causa 2: modelo}\\[2pt]Capacidad variable de modelos compactos para usar el contexto recuperado};
\node[tnaranja] (c3) at (5.5,-3.2) {\textbf{Causa 3: evaluación}\\[2pt]Necesidad de comparar calidad y costo en condiciones controladas};

\node[trojo] (e1) at (-5.5,3.2) {\textbf{Efecto 1}\\[2pt]Recuperación incompleta de información pertinente};
\node[trojo] (e2) at (0,3.2) {\textbf{Efecto 2}\\[2pt]Respuestas que omiten, alteran o no sustentan el contexto};
\node[trojo] (e3) at (5.5,3.2) {\textbf{Efecto 3}\\[2pt]Incertidumbre sobre viabilidad computacional local};

\node[tlbl] at (-8.6,-3.2) {Causas};
\node[tlbl] at (-8.6,3.2) {Efectos};

\foreach \n in {c1,c2,c3} \draw[techa] (\n) -- (centro);
\foreach \n in {e1,e2,e3} \draw[techa] (centro) -- (\n);
\end{tikzpicture}}
\end{figure}

\subsection{Formulación del problema}
\paragraph{Problema general.}
¿De qué manera el diseño e implementación de una arquitectura SLM--RAG con QLoRA influye en la generación contextualizada en Quechua Collao dentro de un entorno computacional offline?

\paragraph{Problemas específicos.}
\begin{itemize}
	\item \textbf{PE1 (Datos y Recuperación):} ¿Cómo estructurar e indexar semánticamente una base de conocimiento documental bilingüe para optimizar la recuperación local en Quechua Collao?
	\item \textbf{PE2 (Núcleo Técnico y Adaptación):} ¿De qué manera los adaptadores QLoRA en un SLM cuantizado optimizan la utilización del contexto recuperado bajo restricciones de hardware local?
	\item \textbf{PE3 (Evaluación Computacional):} ¿Qué nivel de fidelidad léxica, consistencia factual y eficiencia computacional alcanza la arquitectura integrada frente al modelo base y a la recuperación estándar?
\end{itemize}

\section{Objetivos de la investigación}

\subsection{Objetivo general}
Diseñar e implementar una arquitectura SLM--RAG con QLoRA para la generación contextualizada en Quechua Collao dentro de un entorno computacional offline.

\subsection{Objetivos específicos}
\begin{itemize}
	\item \textbf{OE1 (Datos y Recuperación):} Curar, estructurar e indexar semánticamente una base de conocimiento documental bilingüe para optimizar la recuperación local en Quechua Collao.
	\item \textbf{OE2 (Núcleo Técnico y Adaptación):} Implementar y acoplar adaptadores QLoRA en el SLM cuantizado para maximizar la asimilación del contexto RAG en inferencia offline.
	\item \textbf{OE3 (Evaluación Computacional):} Evaluar computacionalmente el rendimiento técnico, la fidelidad lingüística y el anclaje factual de la arquitectura frente al SLM base y al RAG estándar mediante métricas objetivas de PLN y RAG.
\end{itemize}

La Figura~\ref{fig:arbol_objetivos} presenta el árbol de objetivos, organizado en correspondencia directa con el árbol de problemas: los objetivos específicos actúan como hitos cronológicos de ingeniería para lograr el objetivo central, alineado con los Objetivos de Desarrollo Sostenible (ODS).

\begin{figure}[H]
\centering
\caption{Árbol de objetivos: medios de investigación y fines previstos}
\label{fig:arbol_objetivos}
\resizebox{\textwidth}{!}{%
\begin{tikzpicture}[every node/.style={text width=3.3cm, minimum height=1.7cm}]
\node[tcentral, text width=6.6cm, minimum height=1.8cm] (centro) at (0,0) {\textbf{Objetivo Central}\\[2pt]Diseñar e implementar una arquitectura SLM--RAG con QLoRA para generación contextualizada en Quechua Collao};

\node[tteal] (m1) at (-5.5,-3.2) {\textbf{OE1 (Datos)}\\[2pt]Curar, estructurar e indexar la base documental bilingüe};
\node[tteal] (m2) at (0,-3.2) {\textbf{OE2 (Núcleo)}\\[2pt]Acoplar adaptadores QLoRA en el SLM cuantizado};
\node[tteal] (m3) at (5.5,-3.2) {\textbf{OE3 (Evaluación)}\\[2pt]Evaluar métricas computacionales objetivas (PLN y RAGAS)};

\node[tverde, text width=4.0cm] (f1) at (-3.5,3.2) {\textbf{Fin ODS 9 (Innovación)}\\[2pt]Habilitar procesamiento local eficiente en hardware modesto ($\le$ 4~GB)};
\node[tverde, text width=4.0cm] (f2) at (3.5,3.2) {\textbf{Fin ODS 10 (Inclusión)}\\[2pt]Reducir la brecha digital y tecnológica para comunidades quechuahablantes};

\node[tlbl] at (-8.6,-3.2) {Medios};
\node[tlbl] at (-8.6,3.2) {Fines};

\foreach \n in {m1,m2,m3} \draw[techa] (\n) -- (centro);
\foreach \n in {f1,f2} \draw[techa] (centro) -- (\n);
\end{tikzpicture}}
\end{figure}

\section{Hipótesis de la investigación}

\subsection{Hipótesis general}
En las condiciones de prueba definidas, la configuración SLM--RAG con QLoRA obtendrá mejores valores de fidelidad al contexto recuperado y calidad lingüística que las configuraciones de referencia, sin exceder los recursos computacionales disponibles.

\subsection{Hipótesis específicas}
\begin{itemize}
\item \textbf{HE1 (Datos y recuperación):} Con las consultas y etiquetas de relevancia definidas, la indexación semántica local recuperará una mayor proporción de fragmentos pertinentes que la búsqueda léxica de referencia.
\item \textbf{HE2 (Núcleo técnico y adaptación):} Si el ajuste QLoRA resulta ejecutable en el entorno previsto, la configuración adaptada mostrará mayor utilización de la evidencia recuperada que el SLM base, manteniendo constantes los documentos, las consultas y los parámetros de generación.
\item \textbf{HE3 (Evaluación computacional):} La configuración integrada obtendrá mejores puntuaciones de calidad lingüística y fidelidad factual que el SLM base y la configuración SLM--RAG sin adaptación, de acuerdo con las métricas y el análisis estadístico establecidos.
\end{itemize}

\section{Justificación de la investigación}
La investigación se justifica por la necesidad de estudiar alternativas de generación que puedan ejecutarse sin conexión y apoyarse en documentos trazables en Quechua Collao. La propuesta se organiza desde cuatro perspectivas de ingeniería:

\begin{itemize}
\item \textbf{Justificación teórica:} Examina la relación entre recuperación documental, modelos compactos y adaptación de bajo rango en una tarea de generación para Quechua Collao. Los antecedentes sobre traducción de quechua muestran que el tipo de contexto puede influir de manera distinta en el resultado, por lo que la propuesta requiere una evaluación específica \parencite{Court2024}.
\item \textbf{Justificación práctica:} El prototipo se diseñará para funcionar localmente y sin conexión durante la inferencia. Su utilidad y calidad deberán medirse; el plan no presupone que las respuestas serán siempre precisas ni que todos los componentes podrán ejecutarse en la GPU local de 4 GB.
\item \textbf{Justificación metodológica:} Propone pruebas computacionales comparativas para medir el comportamiento del modelo base, la recuperación local y la adaptación QLoRA mediante indicadores de calidad y eficiencia.
\item \textbf{Justificación tecnológica y social:} El posible aporte se relaciona con el acceso a tecnologías locales y con la disponibilidad de recursos lingüísticos. Esta relación se plantea como motivación, no como un impacto social que el experimento vaya a medir:
	\begin{itemize}
\item \textbf{ODS 9 — Industria, innovación e infraestructura:} Se vincula con la exploración de una arquitectura de procesamiento local y eficiente.
\item \textbf{ODS 10 — Reducción de las desigualdades:} Se vincula con la investigación de herramientas digitales para una lengua originaria; el proyecto no medirá reducción de brechas sociales.
	\end{itemize}
\end{itemize}

\section{Alcances y limitaciones}

\subsection{Alcances}
El presente estudio abarca el diseño, implementación y evaluación experimental de una arquitectura SLM--RAG con QLoRA para inferencia local offline en Quechua Collao:

\begin{itemize}
\item \textbf{Alcance tecnológico:} Comprende la selección de un SLM abierto de la familia Qwen2.5 como candidato, la evaluación de cuantización e inferencia local y la integración con recuperación documental local. QLoRA se probará como componente condicionado a la viabilidad de memoria; la cuantización, el formato y el motor final se fijarán y documentarán antes de la comparación \parencite{qwen2024,Dettmers2023}.
\item \textbf{Alcance lingüístico:} Se delimita a materiales identificados como Quechua Collao. Cada fuente y consulta conservará una etiqueta de idioma y variedad; no se extrapolarán resultados a otras variedades quechuas.
\item \textbf{Alcance experimental y métrico:} La evaluación comparará calidad lingüística (chrF++ y SacreBLEU), recuperación y fidelidad al contexto (métricas RAG seleccionadas) y eficiencia local (latencia, tokens por segundo y consumo de RAM/VRAM).
	\item \textbf{Alcance disciplinar:} Se ubica estrictamente dentro de la Ciencia de la Computación, la Inteligencia Artificial y la Ingeniería de Sistemas, en las subáreas de Procesamiento de Lenguaje Natural y Sistemas Distribuidos/Edge Computing.
\end{itemize}

\subsection{Limitaciones}
Se reconocen las siguientes limitaciones del estudio:
\begin{itemize}
\item \textbf{Disponibilidad y cobertura de datos:} Los recursos quechuas existentes no garantizan cobertura suficiente de todas las consultas o registros del Quechua Collao; la selección, licencia y trazabilidad del corpus limitarán el alcance de la evaluación \parencite{Orcotoma2025,Duran2026}.
\item \textbf{Variación lingüística y tokenización:} Los resultados no se generalizarán a otras variedades quechuas. La segmentación de subpalabras y sus efectos sobre memoria y calidad se medirán en el modelo elegido, sin asumir una razón fija frente al español.
\item \textbf{Capacidad de hardware local:} La GPU disponible cuenta con 4~GB de VRAM. El ajuste QLoRA de un modelo compacto puede no ser viable en ese dispositivo; se informará cualquier reducción de alcance, uso de CPU u otra condición computacional, sin presentar las alternativas como equivalentes.
	\item \textbf{Modalidad textual:} El sistema procesa y genera exclusivamente texto digital; el procesamiento de habla (reconocimiento y síntesis de voz) queda fuera de la delimitación del presente proyecto.
\end{itemize}

\section{Matriz de consistencia metodológica}
La Tabla~\ref{tab:matriz_consistencia} resume la correspondencia entre los problemas, objetivos, hipótesis, variables e instrumentos propuestos.

\begin{table}[H]
\centering
\caption{Matriz de consistencia metodológica del proyecto de tesis}
\label{tab:matriz_consistencia}
\scriptsize
\begin{tabularx}{\textwidth}{@{}p{2.8cm}p{3.0cm}p{3.0cm}p{2.4cm}X@{}}
\toprule
\rowcolor{tcAzul!12}
\textbf{Problema} & \textbf{Objetivo} & \textbf{Hipótesis} & \textbf{Variable(s)} & \textbf{Instrumento / Métrica} \\
\midrule
\multicolumn{5}{l}{\textbf{Nivel General}} \\
¿De qué manera el diseño e implementación de una arquitectura SLM--RAG con QLoRA influye en la generación contextualizada en Quechua Collao en un entorno offline? & 
Diseñar e implementar una arquitectura SLM--RAG con QLoRA para generación contextualizada en Quechua Collao en un entorno offline. & 
En las condiciones de prueba definidas, SLM--RAG con QLoRA obtendrá mejores valores de fidelidad y calidad lingüística que las configuraciones de referencia, sin exceder los recursos disponibles. & 
\textbf{VI:} Configuración del sistema.\newline
\textbf{VD:} Calidad lingüística, fidelidad y eficiencia. & 
Consultas repetidas por configuración; métricas lingüísticas, de RAG y de rendimiento. \\
\midrule
\multicolumn{5}{l}{\textbf{Nivel Específico (Fases Metodológicas de Ingeniería)}} \\
\textbf{PE1 (Datos y Recuperación):} ¿Cómo estructurar e indexar semánticamente una base documental bilingüe para optimizar la recuperación local en Quechua Collao? & 
\textbf{OE1:} Curar, estructurar e indexar semánticamente una base documental bilingüe para optimizar la recuperación local en Quechua Collao. & 
\textbf{HE1:} La indexación semántica local recuperará una mayor proporción de fragmentos pertinentes que la búsqueda léxica de referencia, según las etiquetas definidas. & 
Precisión de recuperación (P@k) y trazabilidad de fuentes. & 
Índice local, consultas con documentos pertinentes identificados y registro de fragmentos recuperados. \\
\midrule
\textbf{PE2 (Núcleo y Adaptación):} ¿De qué manera los adaptadores QLoRA del SLM cuantizado optimizan la utilización de contexto bajo restricciones offline? & 
\textbf{OE2:} Implementar y acoplar adaptadores QLoRA en el SLM cuantizado para maximizar la utilización del contexto RAG en inferencia offline. & 
\textbf{HE2:} Si el ajuste QLoRA es ejecutable, la configuración adaptada mostrará mayor uso de evidencia que el SLM base, manteniendo constantes los demás parámetros. & 
Uso de evidencia, latencia, tokens por segundo, RAM y VRAM. & 
Bitácora de ejecución y telemetría en el mismo equipo de pruebas. \\
\midrule
\textbf{PE3 (Evaluación Computacional):} ¿Qué nivel de fidelidad léxica, consistencia factual y eficiencia alcanza la arquitectura frente al modelo base y RAG estándar? & 
\textbf{OE3:} Evaluar computacionalmente el rendimiento, fidelidad lingüística y anclaje factual frente al SLM base y RAG estándar. & 
\textbf{HE3:} La configuración integrada obtendrá mejores puntuaciones que SLM y SLM--RAG en las métricas seleccionadas, si puede evaluarse QLoRA. & 
Fidelidad lingüística (chrF++, BLEU) y anclaje RAG (\textit{faithfulness}, \textit{context recall}). & 
Scripts de evaluación, SacreBLEU, métricas RAG seleccionadas y análisis de medidas repetidas. \\
\bottomrule
\end{tabularx}
\end{table}
